% W4i39_ThirdPeak_Analysis.tex
% DFD 17-Agent Research Programme — Iteration 39 (i39)
% Agents 1 (AeST Analysis), 2 (Disformal Third Peak), 3 (Numerical Estimate),
%         4 (Neutrino HDM), 5 (Baryon Acoustic Enhancement)
% Iteration 8/20 of Quantum Gravity Cycle (i32-i51)
% PRIME DIRECTIVE: Solve the CMB Third Peak Crisis
% Date: 2026-03-22

\documentclass[11pt,a4paper]{article}
\usepackage[utf8]{inputenc}
\usepackage{amsmath,amssymb,amsfonts,amsthm}
\usepackage{hyperref}
\usepackage{booktabs}
\usepackage{longtable}
\usepackage{geometry}
\usepackage{xcolor}
\geometry{margin=2.5cm}

\newtheorem{theorem}{Theorem}
\newtheorem{proposition}{Proposition}
\newtheorem{lemma}{Lemma}
\newtheorem{corollary}{Corollary}
\newtheorem{definition}{Definition}
\newtheorem{remark}{Remark}

\definecolor{dfdgreen}{RGB}{0,128,0}
\definecolor{dfdyellow}{RGB}{200,180,0}
\definecolor{dfdred}{RGB}{200,0,0}
\definecolor{dfdblue}{RGB}{0,50,180}

\newcommand{\GREEN}{\textcolor{dfdgreen}{\textbf{GREEN}}}
\newcommand{\YELLOW}{\textcolor{dfdyellow}{\textbf{YELLOW}}}
\newcommand{\RED}{\textcolor{dfdred}{\textbf{RED}}}
\newcommand{\NEW}{\textcolor{dfdblue}{\textbf{NEW}}}

\title{DFD i39: CMB Third Peak Crisis --- Analysis\\[6pt]
\large Agents 1 (AeST), 2 (Disformal), 3 (Numerical), 4 (Neutrino), 5 (BAO)\\[6pt]
\normalsize Iteration 8/20 of Quantum Gravity Cycle (i32--i51)\\
PRIME DIRECTIVE: Solve the CMB Third Peak Crisis}
\author{DFD 17-Agent Research Programme}
\date{2026-03-22}

\begin{document}
\maketitle
\tableofcontents
\newpage

%=============================================================================
\section{Agent 1: AeST Analysis --- How the Vector Field Fixes the Third Peak}
\label{sec:agent1}
%=============================================================================

\subsection{Mandate}

Deep dive into Skordis \& Z{\l}o\'{s}nik's AeST (Aether Scalar Tensor) theory. Determine \emph{exactly} how the vector field fixes the CMB third peak. Assess whether DFD can replicate this mechanism with its existing scalar + disformal sector.

\subsection{The AeST Field Content}

The AeST theory (Skordis \& Z{\l}o\'{s}nik, PRL 127, 161302, 2021; arXiv:2007.00082) contains three gravitational degrees of freedom beyond the metric $g_{\mu\nu}$:

\begin{enumerate}
\item A \textbf{unit-timelike vector field} $A^\mu$ with $g_{\mu\nu}A^\mu A^\nu = -1$ (the ``aether'')
\item A \textbf{scalar field} $\varphi$ (analogous to DFD's $\psi$)
\item The \textbf{metric} $g_{\mu\nu}$ itself
\end{enumerate}

The action is:
\begin{equation}
S_{\rm AeST} = \frac{1}{16\pi G}\int d^4x\,\sqrt{-g}\left[R + \mathcal{L}_A(A^\mu, \nabla_\nu A^\mu) + \mathcal{L}_\varphi(\varphi, \nabla_\mu\varphi, A^\mu)\right] + S_{\rm matter}[g_{\mu\nu}]
\end{equation}

\textbf{Key papers consulted:}
\begin{itemize}
\item Skordis \& Z{\l}o\'{s}nik (2021), PRL 127, 161302 --- foundational AeST paper
\item Skordis \& Z{\l}o\'{s}nik (2024), MNRAS 531, 272 --- quasistatic spherical solutions
\item Blanchet et al.\ (2024), JCAP 2024, 040 --- Relativistic Khronon theory
\item Durakovi\'{c} et al.\ (2025), arXiv:2503.11174 --- mimetic gravity connection
\item Mistele et al.\ (2024), A\&A 676, A100 --- AeST strong lensing
\item Skordis \& Z{\l}o\'{s}nik (2023), arXiv:2307.15126 --- Hamiltonian formalism
\end{itemize}

\subsection{The Vector Field Mechanism for the Third Peak}

\subsubsection{The Problem Without a Vector}

In standard MOND (or any pure scalar MOND), the CMB acoustic oscillation is driven by the photon-baryon plasma oscillating in gravitational potential wells. The potential wells are sourced \emph{only} by baryons (since there is no CDM). The driving term for the $n$-th acoustic peak goes as:
\begin{equation}
\frac{\delta T}{T}\bigg|_n \propto \left(1 + R_b\right)^{-1/4}\cos\left(n\pi\frac{r_s}{r_{\rm dec}}\right)\Phi_n
\end{equation}
where $R_b = 3\rho_b/(4\rho_\gamma)$ is the baryon-photon ratio and $\Phi_n$ is the gravitational potential at scale $k_n$. Without CDM, $\Phi_n$ decays after horizon entry (baryons oscillate, reducing $\Phi$), and the odd peaks (1st, 3rd, 5th) are suppressed relative to even peaks.

\textbf{Quantitatively:} The third peak height in a baryon-only universe is:
\begin{equation}
\mathcal{R}_3 \equiv \frac{C_{\ell_3}^{\rm baryon\text{-}only}}{C_{\ell_3}^{\rm \Lambda CDM}} \approx 0.55\text{--}0.70
\end{equation}
This is the 30--45\% deficit.

\subsubsection{How the AeST Vector Fixes It}

In AeST, the vector field $A^\mu$ has perturbations $\delta A^\mu$ that evolve independently of the photon-baryon plasma. At the linear perturbation level:

\begin{equation}
\delta A^0 = \alpha_A(k,\eta), \quad \delta A^i = \partial^i \beta_A(k,\eta)
\end{equation}

The key insight: the vector field perturbations $\alpha_A, \beta_A$ \textbf{do not oscillate with the baryons}. Instead, they grow (or remain constant) through recombination, providing a \emph{non-oscillating gravitational potential} that mimics CDM.

The effective Poisson equation in AeST becomes:
\begin{equation}
k^2\Phi = -4\pi G a^2\left[\rho_b\delta_b + \rho_\gamma\delta_\gamma + \rho_\nu\delta_\nu + \rho_{\rm eff}^{(A)}\delta_A^{\rm eff}\right]
\end{equation}
where $\rho_{\rm eff}^{(A)}\delta_A^{\rm eff}$ is the effective energy-momentum contribution from the vector field perturbations. This term:

\begin{enumerate}
\item Does \textbf{not} oscillate with the photon-baryon plasma
\item \textbf{Grows} during radiation domination (like CDM perturbations)
\item Provides the \textbf{potential wells} that drive the baryon infall
\item Has amplitude tuned to match $\Omega_{\rm CDM}h^2 \approx 0.12$
\end{enumerate}

\subsubsection{The Three Gravitational Regimes}

Skordis \& Z{\l}o\'{s}nik (2024) showed AeST has three distinct regimes:

\begin{enumerate}
\item \textbf{Newtonian regime} ($g_N \gg a_0$): Standard GR recovery
\item \textbf{MOND regime} ($g_N \ll a_0$): $\mu(x) \to x$ behavior
\item \textbf{Oscillatory regime}: Gravitational potential oscillations absent in standard MOND
\end{enumerate}

On CMB scales ($\ell \sim 200$--$2000$), the vector field is in the cosmological (non-MOND) regime, acting effectively as CDM. On galaxy scales, it transitions to the MOND regime, giving flat rotation curves.

\subsection{Can DFD Replicate This Without a Vector?}

\subsubsection{DFD's Current Field Content}

DFD has:
\begin{itemize}
\item The metric $g_{\mu\nu}$
\item A scalar field $\psi$ with disformal coupling: $\tilde{g}_{\mu\nu} = e^{2\psi}g_{\mu\nu} + D(\chi)\partial_\mu\psi\partial_\nu\psi$
\end{itemize}

The scalar perturbation $\delta\psi$ has one degree of freedom. CDM has one degree of freedom (the density contrast $\delta_c$). A vector field has \emph{three} degrees of freedom ($\alpha_A, \beta_A$, and a transverse mode).

\textbf{Critical question:} Can one scalar DOF do the work of one CDM DOF?

\subsubsection{The Scalar Cannot Decouple From Baryons}

\textbf{This is the core problem.} In DFD, $\psi$ is sourced by the baryon stress-energy:
\begin{equation}
\Box\psi + V'(\psi) = \frac{D_0}{M_P}\,T_b
\end{equation}
where $T_b$ is the baryon trace. At recombination, baryons oscillate, so $T_b$ oscillates, so $\delta\psi$ oscillates \textbf{in phase with baryons}. This means $\delta\psi$ cannot provide a non-oscillating potential well.

In AeST, the vector field decouples from baryons because its equation of motion involves only $\nabla_\mu A^\nu$ terms (kinetic), not direct matter coupling. The scalar $\varphi$ in AeST couples to baryons, but the vector does not --- it is the vector that provides the CDM-like perturbation.

\subsubsection{Disformal Sector: A Possible Loophole?}

The disformal coupling $D(\chi)\partial_\mu\psi\partial_\nu\psi$ introduces a \emph{direction-dependent} modification. In the cosmological background:
\begin{equation}
\chi = \frac{g^{\mu\nu}\partial_\mu\psi\partial_\nu\psi}{\Lambda_\psi^4} = -\frac{\dot{\psi}^2}{\Lambda_\psi^4}
\end{equation}

The disformal part of the effective metric is:
\begin{equation}
\tilde{g}_{00}^{\rm disf} = D(\chi)\dot{\psi}^2
\end{equation}

This breaks spatial isotropy of perturbations at the level of the effective metric, potentially allowing the $\psi$ perturbation to behave differently from a simple scalar. However:

\textbf{HONEST ASSESSMENT:} The disformal coupling still involves $\partial_\mu\psi$, which is sourced by baryons. The disformal factor modifies the \emph{amplitude} and \emph{propagation speed} of $\delta\psi$, but does not decouple it from the baryon oscillation. The perturbation $\delta\psi$ still oscillates.

\subsection{Agent 1 Conclusion}

\begin{center}
\fbox{\parbox{0.9\textwidth}{
\textbf{Finding 1.1:} The AeST vector field fixes the third peak by providing a non-oscillating gravitational potential that mimics CDM. The vector perturbation grows independently of the baryon oscillation.\\[6pt]
\textbf{Finding 1.2:} DFD's scalar field $\psi$ cannot replicate this mechanism because $\psi$ is sourced by baryons and therefore oscillates with them.\\[6pt]
\textbf{Finding 1.3:} The disformal coupling $D(\chi)$ modifies the amplitude and speed of $\delta\psi$ propagation but does not decouple it from baryon oscillations.\\[6pt]
\textbf{Finding 1.4:} Durakovi\'{c} et al.\ (2025) show that any relativistic MOND theory with a unit-timelike vector can be embedded in a conformal/disformal-invariant mimetic gravity framework. This suggests the vector is not merely a technical convenience --- it encodes a physical degree of freedom that pure scalar theories lack.\\[6pt]
\textbf{STATUS:} \RED{} --- DFD's existing scalar sector \emph{cannot} fix the third peak via the AeST mechanism.
}}
\end{center}

\newpage
%=============================================================================
\section{Agent 2: Disformal Third Peak --- Can $D(\chi)$ Modify the CMB?}
\label{sec:agent2}
%=============================================================================

\subsection{Mandate}

Determine whether DFD's disformal factor $D(\chi)$ can modify the CMB power spectrum at $\ell \sim 800$ (third peak). Compute the disformal correction to the photon-baryon oscillation.

\subsection{Key Papers}

\begin{itemize}
\item Zumalac\'{a}rregui et al.\ (2014), JCAP 1405, 038 --- disformal gravity theories
\item Sakstein \& Verner (2015), PRD 92, 123005 --- disformal coupling
\item Brax et al.\ (2016), arXiv:1604.08038 --- disformal transformations on the CMB
\item Bettoni \& Zumalac\'{a}rregui (2015), PRD 91, 104009 --- kinetic mixing
\item Dalang \& Lombriser (2025), arXiv:2602.19576 --- K-essence Boltzmann dynamics
\item Tsujikawa (2014), JCAP 1404, 043 --- frame invariance
\end{itemize}

\subsection{Disformal Modification of the Photon-Baryon System}

\subsubsection{The DFD Disformal Metric}

DFD's effective metric for matter is:
\begin{equation}
\tilde{g}_{\mu\nu} = e^{2\psi}g_{\mu\nu} + D(\chi)\partial_\mu\psi\partial_\nu\psi
\end{equation}
with $D(\chi) = D_0/(1 + \chi)$ and $\chi = g^{\mu\nu}\partial_\mu\psi\partial_\nu\psi/\Lambda_\psi^4$.

On the cosmological background ($\bar{\psi} = \bar{\psi}(\eta)$):
\begin{equation}
\tilde{g}_{00} = -e^{2\bar{\psi}}\left(1 - D_0\frac{\dot{\bar{\psi}}^2/\Lambda_\psi^4}{1 + \dot{\bar{\psi}}^2/\Lambda_\psi^4}\right)a^2
\end{equation}
\begin{equation}
\tilde{g}_{ij} = e^{2\bar{\psi}}a^2\delta_{ij}
\end{equation}

The disformal part only affects the time-time component.

\subsubsection{Modified Photon Propagation}

Brax et al.\ (2016) showed that disformal transformations imprint on the CMB through two channels:

\textbf{Channel 1: Modified sound speed.} The photon-baryon sound speed becomes:
\begin{equation}
c_s^2 = \frac{1}{3(1+R_b)}\cdot\frac{1}{1 + D_{\rm eff}(\eta)}
\end{equation}
where $D_{\rm eff}$ is the effective disformal factor evaluated on the background. For DFD:
\begin{equation}
D_{\rm eff} = D_0\frac{\bar{\chi}}{1 + \bar{\chi}}, \quad \bar{\chi} = \frac{\dot{\bar{\psi}}^2}{\Lambda_\psi^4}
\end{equation}

\textbf{Channel 2: Modified damping scale.} Dalang \& Lombriser (2025) showed in K-essence cosmology that the diffusion damping scale is modified:
\begin{equation}
k_D^{-2} \propto \int \frac{d\eta}{n_e\sigma_T^{\rm eff}a} \cdot \left(\frac{R_b^2}{1+R_b} + \frac{16}{15}\right)
\end{equation}
where $\sigma_T^{\rm eff}$ includes the disformal rescaling. In DFD, the effective Thomson cross-section is:
\begin{equation}
\sigma_T^{\rm eff} = \sigma_T\cdot e^{-2\bar{\psi}}\cdot(1 + D_{\rm eff})^{-1/2}
\end{equation}

\subsubsection{Numerical Estimate for DFD}

At recombination ($z \approx 1100$), DFD's background scalar is small: $\bar{\psi} \sim D_0\Phi_N \sim 0.017 \times 10^{-5} \sim 10^{-7}$. The disformal factor:
\begin{equation}
\bar{\chi}_{\rm rec} = \frac{\dot{\bar{\psi}}^2}{\Lambda_\psi^4} \sim \frac{(H\bar{\psi})^2}{\Lambda_\psi^4}
\end{equation}

With $H_{\rm rec} \sim 10^{-29}$ eV, $\bar{\psi} \sim 10^{-7}$, $\Lambda_\psi \sim 10^{-3}$ eV:
\begin{equation}
\bar{\chi}_{\rm rec} \sim \frac{(10^{-29} \times 10^{-7})^2}{(10^{-3})^4} = \frac{10^{-72}}{10^{-12}} = 10^{-60}
\end{equation}

\textbf{This is catastrophically small.} The disformal correction to the sound speed is:
\begin{equation}
\frac{\delta c_s^2}{c_s^2} \sim D_0\bar{\chi}_{\rm rec} \sim 0.017 \times 10^{-60} \sim 10^{-62}
\end{equation}

\subsection{Can Perturbations Help?}

The perturbation $\delta\chi$ at first order:
\begin{equation}
\delta\chi = \frac{2\dot{\bar{\psi}}\delta\dot{\psi}}{\Lambda_\psi^4}
\end{equation}

The perturbation $\delta\psi$ at recombination is of order $D_0\Phi \sim 10^{-7}$ (matching the Newtonian potential). So:
\begin{equation}
\delta\chi \sim \frac{2 \times H\bar{\psi} \times H\delta\psi}{\Lambda_\psi^4} \sim 10^{-60}
\end{equation}

Still negligible. The disformal sector operates at a completely different energy scale from the CMB.

\subsection{Comparison With K-essence Models}

Dalang \& Lombriser (2025) found significant disformal modifications because in their K-essence model, the scalar field has $\dot{\varphi} \sim H M_P$ (Planck-scale kinetic energy). In DFD, $\dot{\psi} \sim H \times 10^{-7}$, which is $10^{-7}$ times smaller. Since the disformal correction scales as $\dot{\psi}^2$, the DFD correction is $\sim 10^{-14}$ times that of K-essence models.

\subsection{Agent 2 Conclusion}

\begin{center}
\fbox{\parbox{0.9\textwidth}{
\textbf{Finding 2.1:} The disformal factor $D(\chi)$ modifies the photon-baryon sound speed and damping scale, but the modification is $\sim 10^{-60}$ at recombination.\\[6pt]
\textbf{Finding 2.2:} This is because DFD's scalar has $\dot{\psi}/\Lambda_\psi^2 \sim 10^{-30}$ at recombination, making $\bar{\chi} \sim 10^{-60}$.\\[6pt]
\textbf{Finding 2.3:} K-essence models (Dalang \& Lombriser 2025) achieve significant disformal effects because their scalar has Planck-scale kinetic energy. DFD's scalar is $10^{-7}$ times weaker.\\[6pt]
\textbf{STATUS:} \RED{} --- The disformal sector is irrelevant for the CMB third peak. $D(\chi)$ cannot help.
}}
\end{center}

\newpage
%=============================================================================
\section{Agent 3: Numerical Estimate --- Third Peak With Full MOND $G_{\rm eff}$}
\label{sec:agent3}
%=============================================================================

\subsection{Mandate}

Compute the best possible analytic estimate of the third peak height using DFD's full nonlinear $G_{\rm eff}$, including the transition-regime correction.

\subsection{Key Papers}

\begin{itemize}
\item McGaugh (2004), ApJ 611, 26 --- MOND prediction of second peak
\item Skordis et al.\ (2006), PRL 96, 011301 --- TeVeS CMB
\item Angus (2009), MNRAS 394, 527 --- MOND and the third peak
\item Haslbauer et al.\ (2020), MNRAS 499, 2845 --- $\nu$HDM model
\item Planck Collaboration (2020), A\&A 641, A6 --- CMB data
\item PDG Review (2025) --- Cosmic Microwave Background
\end{itemize}

\subsection{The Third Peak in DFD}

\subsubsection{Setup}

The CMB angular power spectrum peak heights depend on:
\begin{equation}
\mathcal{H}_n = \left|(1 + R_b)^{-1/4}\Phi_n\cos(n\pi) + \text{ISW}\right|^2
\end{equation}

For the $n$-th peak at wavenumber $k_n \approx n\pi/r_s$, the gravitational potential $\Phi_n$ is determined by the Poisson equation:
\begin{equation}
k^2\Phi = 4\pi G_{\rm eff}(k,\eta)\,a^2\rho_b\delta_b
\end{equation}

In DFD, the effective gravitational constant depends on the local acceleration:
\begin{equation}
G_{\rm eff} = G_N\cdot\nu\left(\frac{g_N}{a_0}\right)
\end{equation}
where $\nu(y)$ is the DFD interpolating function (inverse of $\mu$):
\begin{equation}
\nu(y) = \frac{1}{2} + \frac{1}{2}\sqrt{1 + \frac{4}{y}}
\end{equation}

\subsubsection{Acceleration at the Third-Peak Scale}

The third peak corresponds to $\ell_3 \approx 810$, wavenumber:
\begin{equation}
k_3 = \frac{\ell_3}{D_A} \approx \frac{810}{14\,\text{Gpc}} \approx 0.06\,\text{Mpc}^{-1}
\end{equation}

The physical scale at recombination:
\begin{equation}
\lambda_3 = \frac{2\pi}{k_3/(1+z_{\rm rec})} \approx \frac{2\pi \times 1101}{0.06}\,\text{Mpc} \approx 115\,\text{kpc (physical)}
\end{equation}

The Newtonian acceleration at this scale from baryon perturbations ($\delta_b \sim 10^{-3}$ at recombination):
\begin{equation}
g_N = \frac{G_N M_b(<\lambda_3)}{(\lambda_3/2)^2} \sim G_N \times 4\pi\bar{\rho}_b\delta_b\lambda_3/3
\end{equation}

With $\bar{\rho}_b(z=1100) = \Omega_b\rho_c(1+z)^3 \approx 0.05 \times 10^{-29} \times 1.3\times10^9 \approx 6.5 \times 10^{-22}\,\text{g/cm}^3$:
\begin{equation}
g_N \sim \frac{4\pi G_N \times 6.5\times10^{-22} \times 10^{-3} \times 3.5\times10^{23}}{3} \sim 6 \times 10^{-11}\,\text{m/s}^2
\end{equation}

Compare with $a_0 = 1.2 \times 10^{-10}$ m/s$^2$:
\begin{equation}
y_3 = \frac{g_N}{a_0} \approx 0.5
\end{equation}

This is the \textbf{transition regime} ($y \sim 1$), where:
\begin{equation}
\nu(0.5) = \frac{1}{2} + \frac{1}{2}\sqrt{1 + 8} = \frac{1}{2} + \frac{3}{2} = 2.0
\end{equation}

So $G_{\rm eff}/G_N \approx 2.0$ at the third-peak scale. Previous estimates used $\sim$1.6; the corrected value is somewhat larger.

\subsubsection{Effect on the Third Peak}

The potential well depth scales as $\Phi \propto G_{\rm eff}$, so the third peak height scales as:
\begin{equation}
\mathcal{H}_3^{\rm DFD} = \mathcal{H}_3^{\rm baryon} \times \left(\frac{G_{\rm eff}}{G_N}\right)^2 = \mathcal{H}_3^{\rm baryon} \times 4.0
\end{equation}

Starting from $\mathcal{R}_3 \approx 0.60$ (baryon-only third peak relative to $\Lambda$CDM):
\begin{equation}
\mathcal{R}_3^{\rm DFD} = 0.60 \times 4.0 = 2.4
\end{equation}

\textbf{WAIT.} This gives a third peak that is \emph{too high} --- 2.4 times the $\Lambda$CDM value. This cannot be right.

\subsubsection{The Error: Self-Consistent Treatment}

The problem: the MOND enhancement applies to \emph{all} peaks, not just the third. The first peak is also enhanced. What matters is the \textbf{peak ratio}:
\begin{equation}
R_{31} \equiv \frac{C_{\ell_3}}{C_{\ell_1}}
\end{equation}

In $\Lambda$CDM: $R_{31}^{\rm \Lambda CDM} \approx 0.43$.

If $G_{\rm eff}$ at the first peak ($\ell_1 \approx 220$, $k_1 \approx 0.016$ Mpc$^{-1}$):
\begin{equation}
y_1 = \frac{g_N(k_1)}{a_0} \sim \frac{g_N(k_3) \times (k_1/k_3)^2}{a_0} \times \frac{\delta_b(k_1)}{\delta_b(k_3)}
\end{equation}

The first peak is at larger scale, so $g_N$ is smaller (deeper in MOND). With $k_1/k_3 \approx 0.27$:
\begin{equation}
y_1 \approx y_3 \times (0.27)^2 \times \frac{\delta_b(k_1)}{\delta_b(k_3)} \approx 0.5 \times 0.073 \times 3 \approx 0.11
\end{equation}

(Factor of $\sim$3 because the baryon transfer function is larger at the first-peak scale.)

\begin{equation}
\nu(0.11) = \frac{1}{2} + \frac{1}{2}\sqrt{1 + 36.4} = \frac{1}{2} + 3.06 = 3.56
\end{equation}

So $G_{\rm eff}/G_N \approx 3.56$ at the first peak. This is \emph{larger} than at the third peak.

The peak ratio:
\begin{equation}
R_{31}^{\rm DFD} = R_{31}^{\rm baryon} \times \frac{\nu(y_3)^2}{\nu(y_1)^2} = R_{31}^{\rm baryon} \times \frac{4.0}{12.7} = R_{31}^{\rm baryon} \times 0.315
\end{equation}

In a baryon-only model, $R_{31}^{\rm baryon} \approx 0.25$. So:
\begin{equation}
R_{31}^{\rm DFD} \approx 0.25 \times 0.315 = 0.079
\end{equation}

\textbf{Observed:} $R_{31}^{\rm obs} \approx 0.43$. DFD predicts $0.079$. This is a factor of \textbf{5.4 too low}.

\subsubsection{Understanding the Result}

The MOND enhancement is \emph{scale-dependent}: it enhances large scales (low $\ell$, deep MOND) more than small scales (high $\ell$, transition regime). This \textbf{worsens} the third-peak deficit rather than helping it.

The reason: at the third-peak scale ($y \sim 0.5$), the enhancement is $\nu \approx 2$. But at the first-peak scale ($y \sim 0.1$), the enhancement is $\nu \approx 3.6$. The first peak gets \emph{more} boost than the third peak. The peak ratio decreases.

\subsection{Agent 3 Conclusion}

\begin{center}
\fbox{\parbox{0.9\textwidth}{
\textbf{Finding 3.1:} The MOND $G_{\rm eff}$ enhancement at the third-peak scale is $\nu(y_3) \approx 2.0$ (transition regime, $y_3 \approx 0.5$).\\[6pt]
\textbf{Finding 3.2 (CRITICAL):} The scale-dependent MOND enhancement \textbf{worsens} the peak ratio $R_{31}$ because the first peak is enhanced more ($\nu(y_1) \approx 3.6$) than the third peak. The predicted ratio $R_{31}^{\rm DFD} \approx 0.079$ vs.\ observed $\approx 0.43$.\\[6pt]
\textbf{Finding 3.3:} This is a factor of 5.4 discrepancy, worse than the naive 30--45\% deficit. The MOND enhancement makes the problem \emph{worse}, not better.\\[6pt]
\textbf{Finding 3.4:} The previous estimate of ``30--45\% deficit'' was computed without accounting for the differential MOND enhancement. The true deficit in the peak ratio is a factor of $\sim$5.\\[6pt]
\textbf{CAUTION:} This is an analytic order-of-magnitude estimate. The actual computation requires a full Boltzmann solver with DFD's $G_{\rm eff}(k,\eta)$. But the \emph{direction} of the effect is robust: scale-dependent MOND enhancement suppresses the third peak relative to the first.\\[6pt]
\textbf{STATUS:} \RED{} --- MOND $G_{\rm eff}$ makes the third-peak problem worse, not better.
}}
\end{center}

\newpage
%=============================================================================
\section{Agent 4: Neutrino Hot Dark Matter --- Can 2 eV Neutrinos Help?}
\label{sec:agent4}
%=============================================================================

\subsection{Mandate}

Assess whether DFD's predicted 2 eV neutrinos (prediction P35) can partially fill the third-peak deficit. Re-analyze in light of Katz et al.\ constraints and the Russell et al.\ (2026) $\nu$HDM simulations.

\subsection{Key Papers}

\begin{itemize}
\item Haslbauer et al.\ (2020), MNRAS 499, 2845 --- $\nu$HDM with 11 eV sterile neutrino
\item Russell et al.\ (2026), MNRAS submitted, arXiv:2602.21975 --- $\nu$HDM $N$-body simulations
\item Angus (2009), MNRAS 394, 527 --- light sterile neutrinos in MOND
\item DFD prediction P35: $m_\nu \approx 2$ eV (sum of active neutrinos)
\item Planck (2020): $\sum m_\nu < 0.12$ eV (95\% CL, $\Lambda$CDM)
\item KATRIN (2025): $m_\nu < 0.45$ eV (90\% CL, direct)
\end{itemize}

\subsection{The $\nu$HDM Paradigm and Its Failure}

\subsubsection{The Idea}

Replace CDM with hot dark matter (massive neutrinos). In MOND, the enhanced gravitational growth partially compensates for the free-streaming erasure of small-scale structure. The CMB can be fitted by adjusting the neutrino mass and species.

The $\nu$HDM model (Angus 2009, Haslbauer et al.\ 2020) uses:
\begin{itemize}
\item 11 eV sterile neutrino (provides HDM density $\Omega_\nu h^2 \sim 0.12$)
\item MOND gravity (enhances growth)
\item Standard expansion history
\end{itemize}

\subsubsection{Russell et al.\ (2026): Definitive Failure}

Russell et al.\ (2026) performed the largest $\nu$HDM $N$-body simulations (256$^3$ particles, $800/h$ Mpc box). Key results:

\begin{enumerate}
\item \textbf{Massive overproduction of structure:} Most massive cluster $\sim 5 \times 10^{17}\,M_\odot/h$ (observed: $\sim 10^{15}\,M_\odot/h$) --- 500$\times$ too massive.
\item \textbf{Excessive peculiar velocities:} Several thousand km/s (observed: $\sim 600$ km/s) --- 5$\times$ too fast.
\item \textbf{Ruled out at $>5\sigma$} from bulk flow alone.
\item \textbf{Top-heavy mass function:} Too many massive clusters, too few low-mass ones.
\end{enumerate}

\textbf{Conclusion from Russell et al.:} ``Observations clearly require a much milder enhancement to the rate of structure growth in $\Lambda$CDM than is provided by the $\nu$HDM paradigm.''

\subsection{DFD's 2 eV Neutrinos: A Different Beast}

DFD predicts $\sum m_\nu \approx 2$ eV, not 11 eV. The neutrino energy density:
\begin{equation}
\Omega_\nu h^2 = \frac{\sum m_\nu}{94\,\text{eV}} = \frac{2}{94} \approx 0.021
\end{equation}

Compare with $\Omega_{\rm CDM}h^2 \approx 0.12$. The DFD neutrinos provide only $\sim$18\% of the CDM energy density.

\subsubsection{Effect on the Third Peak}

Massive neutrinos contribute to the gravitational potential at scales larger than their free-streaming length:
\begin{equation}
k_{\rm fs} = \frac{0.018\,\Omega_m^{1/2}}{(m_\nu/1\,\text{eV})}\,h\,\text{Mpc}^{-1} \approx 0.003\,h\,\text{Mpc}^{-1} \quad (m_\nu = 2\,\text{eV})
\end{equation}

The third peak is at $k_3 \approx 0.06$ Mpc$^{-1}$. Since $k_3 \gg k_{\rm fs}$, the 2 eV neutrinos have \textbf{completely free-streamed} at the third-peak scale. They contribute \textbf{nothing} to the gravitational potential at $\ell \sim 800$.

\textbf{Even at scales where they contribute ($k < k_{\rm fs}$):}
\begin{equation}
\frac{\delta\Phi_\nu}{\delta\Phi_{\rm CDM}} \approx \frac{\Omega_\nu}{\Omega_{\rm CDM}} \approx 0.18
\end{equation}

And this 18\% is only relevant for $\ell \lesssim 50$, far from the third peak.

\subsubsection{The Russell et al.\ Problem Applied to DFD}

If DFD's MOND-enhanced growth combined with even 2 eV neutrinos produces structure overproduction at $z = 0$, DFD faces the same problem as $\nu$HDM (just with smaller neutrino mass). The nonlinear growth $\delta \propto a^2$ in MOND means:

At $z = 0$, the growth factor is $(1+z_{\rm eq})^2 \sim 10^6$ times the $\Lambda$CDM value. Even with $\Omega_\nu/\Omega_{\rm CDM} = 0.18$, the total growth is $\sim 10^5$ times too large.

\textbf{HOWEVER:} DFD's MOND regime only applies where $g < a_0$. On Gpc scales, $g \gg a_0$ and DFD reduces to GR. So the structure overproduction problem may not apply to DFD at the same level as $\nu$HDM with global MOND.

\subsection{Agent 4 Conclusion}

\begin{center}
\fbox{\parbox{0.9\textwidth}{
\textbf{Finding 4.1:} DFD's 2 eV neutrinos have completely free-streamed at the third-peak scale ($k_3 \gg k_{\rm fs}$). They contribute \textbf{zero} to the third-peak gravitational potential.\\[6pt]
\textbf{Finding 4.2:} Even at low $\ell$, the neutrino contribution is only $\sim$18\% of $\Omega_{\rm CDM}$ --- grossly insufficient.\\[6pt]
\textbf{Finding 4.3:} Russell et al.\ (2026) ruled out $\nu$HDM at $>5\sigma$ due to massive structure overproduction. DFD must ensure its MOND regime does not produce similar overproduction.\\[6pt]
\textbf{Finding 4.4:} DFD's scale-dependent $G_{\rm eff}$ (MOND only where $g < a_0$) may protect it from the Russell et al.\ problem on Gpc scales, but this requires explicit numerical verification.\\[6pt]
\textbf{STATUS:} \RED{} --- Neutrino HDM cannot fix the third peak. Free-streaming kills the mechanism.
}}
\end{center}

\newpage
%=============================================================================
\section{Agent 5: Baryon Acoustic Enhancement --- MOND and BAO}
\label{sec:agent5}
%=============================================================================

\subsection{Mandate}

Determine whether DFD's MOND enhancement changes the baryon acoustic oscillation amplitude in a way that affects the peak ratio.

\subsection{Key Papers}

\begin{itemize}
\item Eisenstein \& Hu (1998), ApJ 496, 605 --- BAO fitting formulae
\item DESI Collaboration (2025), multiple papers --- BAO measurements
\item Bruni et al.\ (2024), arXiv:2407.15832 --- nonminimally coupled gravity and DESI BAO
\item Sakstein et al.\ (2025), arXiv:2501.15298 --- scalar-tensor gravity and DESI BAO
\item Hu \& Sugiyama (1996), ApJ 471, 542 --- driving effect
\item McGaugh (2004), ApJ 611, 26 --- MOND acoustic peaks
\end{itemize}

\subsection{BAO Physics in DFD}

\subsubsection{The Sound Horizon}

The sound horizon at recombination determines the acoustic peak spacing:
\begin{equation}
r_s = \int_0^{a_{\rm rec}} \frac{c_s\,da}{a^2 H(a)}
\end{equation}

In DFD, $H(a)$ is identical to $\Lambda$CDM (DFD does not modify the expansion history --- it has bare $\Lambda$). The sound speed $c_s = 1/\sqrt{3(1+R_b)}$ is unchanged (Agent 2 showed the disformal correction is $\sim 10^{-60}$). Therefore:
\begin{equation}
r_s^{\rm DFD} = r_s^{\rm \Lambda CDM} \approx 147\,\text{Mpc}
\end{equation}

The peak spacing is unmodified. \textbf{This is good} --- DFD does not conflict with DESI BAO data.

\subsubsection{The Baryon Loading Effect}

The baryon loading parameter $R_b = 3\rho_b/(4\rho_\gamma)$ controls the asymmetry between compression (odd) and rarefaction (even) peaks. In $\Lambda$CDM, CDM provides additional compression, boosting odd peaks.

In DFD, there is no CDM. The MOND enhancement partially substitutes:
\begin{equation}
\Phi_{\rm eff}(k,\eta) = \nu(y(k,\eta))\cdot\Phi_N(k,\eta)
\end{equation}

But as Agent 3 showed, $\nu$ is scale-dependent and enhances the first peak more than the third. The baryon loading asymmetry is:
\begin{equation}
\frac{C_{\ell_3}}{C_{\ell_2}} \propto \frac{(1+R_b)^{1/2} + \nu(y_3)\Phi_3}{(1+R_b)^{1/2} - \nu(y_2)\Phi_2}
\end{equation}

Since $\nu(y_2) > \nu(y_3)$ (the second peak is at larger scale, deeper in MOND), the rarefaction peak (even) gets more MOND boost than the compression peak (odd). This \textbf{inverts} the usual CDM-like asymmetry.

\subsubsection{Peak Ratio $R_{32}$}

In $\Lambda$CDM: $R_{32} \equiv C_{\ell_3}/C_{\ell_2} \approx 0.92$.

In DFD, the MOND-inverted asymmetry gives:
\begin{equation}
R_{32}^{\rm DFD} \approx \frac{\nu(y_3)^2}{\nu(y_2)^2} \times R_{32}^{\rm baryon} \approx \frac{4.0}{6.25} \times 0.55 \approx 0.35
\end{equation}

(Using $\nu(y_2) \approx 2.5$ for the second peak at $y_2 \approx 0.25$.)

Observed: $R_{32} \approx 0.92$. DFD predicts $\approx 0.35$. This is off by a factor of $\sim$2.6.

\subsection{DESI BAO Implications}

DESI DR2 (2025) found mild preference for dynamical dark energy ($w_0 w_a$CDM over $\Lambda$CDM). Bruni et al.\ (2024) showed this can be interpreted as nonminimally coupled gravity in the Horndeski class.

DFD's scalar-tensor structure falls within the Horndeski class (scalar with nonminimal coupling and disformal terms). The DESI BAO data are \emph{consistent} with DFD's background expansion, but the perturbation sector (peak ratios) remains problematic.

\subsection{Agent 5 Conclusion}

\begin{center}
\fbox{\parbox{0.9\textwidth}{
\textbf{Finding 5.1:} DFD's sound horizon $r_s$ is identical to $\Lambda$CDM (147 Mpc). The peak spacing is unmodified.\\[6pt]
\textbf{Finding 5.2:} The MOND enhancement inverts the usual compression/rarefaction asymmetry because even peaks (at larger scales) get more MOND boost than odd peaks.\\[6pt]
\textbf{Finding 5.3:} The peak ratio $R_{32}^{\rm DFD} \approx 0.35$ vs.\ observed $\approx 0.92$. This is a factor of 2.6 discrepancy.\\[6pt]
\textbf{Finding 5.4:} DFD is consistent with DESI BAO background data but fails at the perturbation level.\\[6pt]
\textbf{STATUS:} \RED{} --- BAO analysis confirms and deepens the third-peak crisis. The scale-dependent MOND enhancement damages ALL peak ratios.
}}
\end{center}

\newpage
%=============================================================================
\section{Summary of Agents 1--5}
\label{sec:summary}
%=============================================================================

\begin{center}
\begin{tabular}{llcl}
\toprule
\textbf{Agent} & \textbf{Mechanism} & \textbf{Status} & \textbf{Key Finding} \\
\midrule
1 (AeST) & Vector field decoupled from baryons & \RED & DFD scalar couples to baryons; cannot replicate \\
2 (Disformal) & $D(\chi)$ modifies sound speed & \RED & Correction is $\sim 10^{-60}$; utterly negligible \\
3 (Numerical) & MOND $G_{\rm eff}$ enhancement & \RED & Makes problem \emph{worse} (factor of 5.4 in $R_{31}$) \\
4 (Neutrino) & 2 eV $\nu$HDM & \RED & Free-streamed at third-peak scale; zero contribution \\
5 (BAO) & Peak ratio asymmetry & \RED & MOND inverts compression/rarefaction; $R_{32}$ off by 2.6$\times$ \\
\bottomrule
\end{tabular}
\end{center}

\textbf{Overall assessment:} All five mechanisms investigated fail to resolve the CMB third-peak deficit. More critically, Agents 3 and 5 reveal that the problem is \textbf{worse than previously estimated}: the scale-dependent MOND enhancement actively suppresses the third peak relative to the first. The deficit is not 30--45\% but closer to a \textbf{factor of 5} in the peak ratio $R_{31}$.

\end{document}
